\documentclass[11pt]{article}

% ── Encoding & fonts ──────────────────────────────────────────────────────────
\usepackage[T1]{fontenc}
\usepackage[utf8]{inputenc}
\usepackage{lmodern}
\usepackage{microtype}

% ── Page geometry ─────────────────────────────────────────────────────────────
\usepackage[margin=1.25in]{geometry}

% ── Mathematics ───────────────────────────────────────────────────────────────
\usepackage{amsmath}
\usepackage{amssymb}
\usepackage{amsthm}

% ── Tables ────────────────────────────────────────────────────────────────────
\usepackage{booktabs}
\usepackage{array}
\usepackage{tabularx}

% ── Lists ─────────────────────────────────────────────────────────────────────
\usepackage{enumitem}
\setlist{itemsep=2pt, topsep=4pt}

% ── Hyperlinks ────────────────────────────────────────────────────────────────
\usepackage[colorlinks=true, linkcolor=black, citecolor=black,
            urlcolor=black, pdfborder={0 0 0}]{hyperref}
\usepackage{url}

% ── Bibliography ──────────────────────────────────────────────────────────────
\usepackage[numbers, sort&compress]{natbib}
\bibliographystyle{unsrtnat}

% ── Paragraph spacing ─────────────────────────────────────────────────────────
\usepackage{parskip}
\setlength{\parskip}{6pt}

% ── Section formatting ────────────────────────────────────────────────────────
\usepackage{titlesec}
\titleformat{\section}{\large\bfseries}{\thesection.}{0.6em}{}
\titleformat{\subsection}{\normalsize\bfseries}{\thesubsection}{0.6em}{}
\titlespacing{\section}{0pt}{14pt}{6pt}
\titlespacing{\subsection}{0pt}{10pt}{4pt}

% ── Abstract ──────────────────────────────────────────────────────────────────
\usepackage{abstract}
\renewcommand{\abstractnamefont}{\normalfont\bfseries}
\renewcommand{\abstracttextfont}{\normalfont\small}

% ── Footnotes ─────────────────────────────────────────────────────────────────
\usepackage[hang, flushmargin]{footmisc}

% ── Theorem-like environments ─────────────────────────────────────────────────
\theoremstyle{definition}
\newtheorem{definition}{Definition}

% ── The information-rights diagnostic symbol ─────────────────────────────────
\newcommand{\Iax}{\ensuremath{\hat{I}}}

% ─────────────────────────────────────────────────────────────────────────────
\title{\textbf{Monetary Ecology:\\
An MCB/P/S/\Iax{} Framework for Evaluating\\
Monetary Architectures}}

\author{%
  Dr.\ Richard Nicholson\\[4pt]
  \small Tathata Systems Ltd, United Kingdom\\
  \small \href{mailto:tathatasystems@proton.me}{tathatasystems@proton.me}
}

\date{2026\\[4pt]
\small \textbf{Working draft.} Formalised from the essay
\emph{From Blood Money to Monetary Ecology} (2025).}
% ─────────────────────────────────────────────────────────────────────────────

\begin{document}

\maketitle

\begin{abstract}
A monetary architecture encodes, in its rules of issuance, custody, and redemption,
three commitments that conventional metrics---stability, throughput, settlement
finality---leave implicit: who is counted, who decides, and who claims the surplus.
This paper presents a framework for making those commitments explicit and comparable.
It scores a monetary architecture on three axes---\emph{Moral Circle Breadth} (MCB),
\emph{Polycentricity} (P), and \emph{Surplus-Claim Regime} (S)---together with an
\emph{Information-Rights} diagnostic ($\Iax$) that qualifies the others by asking
whether their claims are verifiable without exposing persons to surveillance. We give
a 0--5 scoring method with a fully worked example, and apply the framework to a set of
architectural ideal types. We argue one substantive structural claim, that low
polycentricity exerts pressure toward concentrated surplus-claim ($P \rightsquigarrow
S$), while keeping the axes independent as measurements. This is a framework paper:
the scores are illustrative, demonstrating the method rather than measuring real
systems, and empirical scoring of deployed monetary systems is left to future work.
\end{abstract}

% ─────────────────────────────────────────────────────────────────────────────
\section{Introduction}
\label{sec:intro}

A monetary system is usually evaluated on narrow technical criteria. These matter,
but they are silent on the questions a monetary architecture in fact answers in its
design. A system can be perfectly stable and perfectly extractive; it can be flawless by
every technical measure and answer to no one. The reason these metrics mislead is one
of timescale: a monetary architecture's consequences unfold not at the speed of a
transaction but at the pace of a \emph{metabolism}---accumulation, exclusion, and the
slow concentration of claim, over months and years---so a system can excel at
everything measured in milliseconds and fail at everything that takes years to
surface. It is those slow dimensions that this framework measures.

The framework of this paper rests on a single observation, which we take as canonical:
\emph{a monetary architecture encodes, in its rules of issuance, custody, and
redemption, who is counted, who decides, and who claims the surplus it generates.}
Money is not neutral~\cite{ingham2004}. The framework's purpose is to make those three commitments---and
the verifiability of a system's claims about them---explicit, comparable, and
contestable on a common scale.

This is a \emph{framework} paper. Its contribution is an evaluative instrument: four
well-defined dimensions, a scoring method, and a worked demonstration that the method
applies end to end. The scores it reports are \emph{illustrative}---they exercise the
instrument on architectural ideal types, not deployed systems; we claim the instrument
and one structural tendency that follows from it, not measurements. Empirical scoring
of real systems is a programme in its own right, scoped in \S\ref{sec:limits}.

\paragraph{Contributions.} This paper (i) defines the MCB/P/S framework and the
$\Iax$ diagnostic (\S\ref{sec:framework}); (ii) gives a 0--5 scoring method with a
fully worked example, foregrounding the illustrative status of all scores
(\S\ref{sec:scoring}); (iii) argues the structural tendency $P \rightsquigarrow S$ as
an interpretation that emerges from the framework (\S\ref{sec:causality}); (iv)
applies the framework to a set of architectural ideal types, with a federated
``monetary ecology'' presented as a hypothesis rather than a demonstrated winner
(\S\ref{sec:application}); and (v) states its scope and the work it defers
(\S\ref{sec:limits}). A short closing note relates the framework to coordinator-free
substrate design (\S\ref{sec:broader}).

A remark on terminology. The originating essay uses the deliberately polemical term
\emph{blood money} for the limit case of a centralised medium that finances outcomes
its holders would not individually endorse. We retain it only as a rhetorical
limit-case label, not as an analytic category; the framework's work is done by the
four dimensions below.

% ─────────────────────────────────────────────────────────────────────────────
\section{The Framework: MCB, P, S, and \texorpdfstring{$\Iax$}{I-hat}}
\label{sec:framework}

The framework rests on three questions a monetary architecture cannot avoid
answering, and which conventional metrics do not ask, plus a diagnostic that gates
the answers.

\begin{definition}[Moral Circle Breadth (MCB) --- ``who counts?'']
The effective scope of those granted protection, voice, and material consideration by
the system. A broad moral circle extends standing beyond present participants to
non-citizens, future generations, non-human animals, and ecosystems.
\end{definition}

MCB has a load-bearing \emph{voiceless edge}: the parties whose standing most tests a
moral circle---future generations, non-human animals, ecosystems---cannot advocate
for themselves within the system, so it extends them standing only structurally,
through rules rather than participation. MCB therefore measures who is
\emph{considered}, not who is represented, and its operational content is what a
high-MCB architecture would actually \emph{do}: ecological reserve requirements,
intergenerational provisioning, externality pricing in credit allocation,
inclusion-of-access rights for non-citizens. Of the four dimensions it is the most
explicitly normative and the least mechanical, by nature rather than by omission.

\begin{definition}[Polycentricity (P) --- ``who decides?'']
The dispersion, contestability, and recallability of decision rights over the system
across multiple independent centres, rather than their concentration in a single
authority. After Ostrom~\cite{ostrom1990}, what matters is not the mere existence of
many centres but that their authority is contestable and recallable.
\end{definition}

\begin{definition}[Surplus-Claim Regime (S) --- ``who benefits?'']
The incidence of residual income and rents---who captures the value a system
generates after costs. A distributed regime returns value to participants,
communities, and the commons; a concentrated one routes it to the owners of the
issuing or custodial layer.
\end{definition}

The three axes are \emph{orthogonal as measurements}: improving one does not
automatically improve another, and a system may score high on one while low on
another (broad protections can coexist with pervasive surveillance; a widely held
shareholder base can still concentrate surplus-claim). It is precisely because the
axes can diverge that measuring them separately is informative; the framework's
descriptive content is the \emph{vector} $(\text{MCB}, P, S)$, not its average. Each
axis draws on an established tradition---the moral-circle literature after
Singer~\cite{singer1981} for MCB, Ostrom's~\cite{ostrom1990} polycentric governance
for P, the political economy of surplus and rent for S---but the synthesis into a
single evaluative instrument for monetary architecture is, to our knowledge, novel.

\begin{definition}[Information-Rights diagnostic ($\Iax$)]
A red-flag diagnostic, not a fourth axis, that qualifies the confidence with which
MCB/P/S can be asserted. Its principle is \emph{privacy for people, transparency for
power}: persons should be shielded from surveillance, while the mechanisms that
exercise power over the system should be continuously auditable.
\end{definition}

Making verifiability a \emph{diagnostic} rather than an axis is deliberate: it does
not describe a distributional commitment, it prevents good distributional claims from
being laundered through opaque or capturable means. A high MCB/P/S vector carried by a
low $\Iax$ is provisional, because the means by which it is achieved cannot be
verified and may be captured without notice. $\Iax$ scores five domains---proofs of
reserves and flows; portability and exit; custody topology; governance traceability;
surveillance exposure---net of penalties for specific opacity failures (no public
proofs; a single custodian able to override; failed portability; audit-blocking
clauses; unreported material incidents), and is clamped to $[0,5]$:
\[
  \Iax = \max\!\bigl(0,\ \min(5,\ \Iax_{\text{raw}} - \textstyle\sum\text{penalties})\bigr).
\]
It is reported \emph{alongside} the vector as a confidence qualifier. The framework's
hardest design questions live here---whether a small co-operative is ``people'' or
``power'', how to prove reserves without exposing individual transaction graphs---and
we flag them as open (\S\ref{sec:limits}) rather than resolved.

% ─────────────────────────────────────────────────────────────────────────────
\section{Scoring: Method, Confidence, and a Worked Example}
\label{sec:scoring}

\paragraph{Method.} Each axis is scored on a 0--5 scale as a weighted average of two
to five observable indicators, each normalised to the scale. Indicator sets are
domain-specific and \emph{published with the score}, so that a placement can be
contested at the level of its evidence rather than its conclusion---an application of
the framework's own polycentricity to itself. Systems are plotted on the plane
$(x = P,\ y = S)$ with bubble magnitude given by MCB and $\Iax$ reported separately
as a confidence-qualified annotation; the $(P,S)$ plane is chosen because it is the
plane in which the structural tendency of \S\ref{sec:causality} runs.

\paragraph{A worked example.} Table~\ref{tab:worked} scores one ideal type, a
\emph{closed custodial platform}~\cite{srnicek2017} (custodial wallets, closed rails, a single corporate
operator), end to end, on indicators uncontroversial enough that the example
demonstrates the \emph{method} without a dispute over the \emph{evidence}. Like all
scores in this paper it is illustrative; its point is the \emph{reproducible scoring
form}---a reader can follow exactly how a score is built and dispute any indicator,
whether or not they accept the number. We report axis scores to the nearest
half-point: the method is exact, but the interpretation it supports is coarse.

\begin{table}[!ht]
\centering\small
\renewcommand{\arraystretch}{1.25}
\begin{tabularx}{\textwidth}{@{}p{1.6cm}p{4.4cm}Xc@{}}
\toprule
\textbf{Axis (score)} & \textbf{Indicator (weight)} & \textbf{Evidence in a closed custodial platform} & \textbf{0--5} \\
\midrule
\textbf{P} \newline ($\approx$0.5)
  & Recallability of authority (0.40) & Operator sets and changes terms unilaterally; participants cannot recall it & 0 \\
  & Dispersion of decision rights (0.35) & A single corporate centre governs issuance, custody, and rules & 0 \\
  & Participant contestability (0.25) & Complaint channels and external regulatory recourse, but no power over rule-making & 2 \\
\addlinespace
\textbf{S} \newline ($\approx$0.5)
  & Residual-claim incidence (0.50) & Fees, spread, and data rents accrue to platform shareholders & 1 \\
  & User share of value generated (0.30) & Users generate network value but capture little of it & 0.5 \\
  & Anti-rent provisions (0.20) & Closed rails, lock-in; no interoperability or fee caps & 0 \\
\addlinespace
\textbf{MCB} \newline ($\approx$1.5)
  & Access breadth (0.40) & Broad consumer access where profitable; the unprofitable are excluded & 2.5 \\
  & Standing for non-present stakeholders (0.30) & No provision for future generations or ecological cost & 0 \\
  & Non-discrimination (0.30) & KYC-gated; commercial exclusion permitted & 1 \\
\addlinespace
\textbf{$\Iax$} \newline (0)
  & Portability \& exit (0.25) & Account data exportable, but value locked into the platform & 3 \\
  & Governance traceability (0.20) & Published terms and change-logs; no independent audit & 2.5 \\
  & Proofs, custody, surveillance (0.55) & No public proofs of reserves; a single custodian; full transaction surveillance & 0 \\
  & \emph{Penalties} & No public proofs ($-1.0$) and single-custodian override ($-1.0$) & $-2.0$ \\
\bottomrule
\end{tabularx}
\caption{A worked \emph{illustrative} score for one architectural ideal type, with
axis scores rounded to the nearest half-point (the method is exact; the
interpretation is coarse). Weighted averages: $P = 0(0.40)+0(0.35)+2(0.25) = 0.5$;
$S = 1(0.50)+0.5(0.30)+0(0.20) = 0.65 \approx 0.5$; $\text{MCB} =
2.5(0.40)+0(0.30)+1(0.30) = 1.3 \approx 1.5$. The diagnostic earns partial raw credit
for exportable data and published terms ($\Iax_{\text{raw}} = 3(0.25)+2.5(0.20) =
1.25$) but two fatal red flags---no public proofs and a single custodian able to
override---deduct $2.0$, giving $\Iax = \max(0, 1.25-2.0) = 0$. The evidence column is
what a reader contests: a placement is disputed by disputing an indicator's evidence
or weight, not by rejecting the conclusion.}
\label{tab:worked}
\end{table}

The example makes two things concrete. The vector $(\text{MCB} \approx 1.5,\ P
\approx 0.5,\ S \approx 0.5)$ is valuable not as a measurement but because it is
\emph{traceable} to indicators a reader can argue with. And the $\Iax$ row shows the
penalty mechanism doing its work: partial verifiability (raw $1.25$) does not rescue a
system that carries fatal red flags, so the platform's distributional scores are
flagged as unverifiable in any case. Penalties matter precisely where a system has
\emph{some} apparent transparency but a disqualifying opacity beneath it.

% ─────────────────────────────────────────────────────────────────────────────
\section{A Structural Tendency: \texorpdfstring{$P \rightsquigarrow S$}{P pressures S}}
\label{sec:causality}

The axes are scored separately because they can diverge at a point in time. But
divergence at a point in time does not mean the absence of pressure over time. We
argue one such pressure, written $P \rightsquigarrow S$ to mark it as a directional
tendency rather than a law:

\begin{quote}
Low polycentricity exerts pressure toward concentrated surplus-claim. Where decision
rights over a monetary system concentrate, and absent explicit anti-rent rules, the
control points that concentration creates tend to be converted into rent.
\end{quote}

The mechanism is the one Olson~\cite{olson1965} identified: the benefits of
controlling a coordination mechanism are concentrated on the controller while the
costs are diffuse, so concentrated power is under standing pressure to convert itself
into concentrated claim. The conversion is not mysterious; it runs through a short
causal chain. Concentrated decision rights create \emph{control points}---over
issuance, custody, permissioning, and settlement. Control points create
\emph{exclusion or dependency}: a participant must transact through the controller or
be shut out. Exclusion and dependency create \emph{pricing power}---seigniorage at
issuance, fee extraction and lock-in at custody, exclusionary rents at permissioning,
chokepoint pricing at settlement. And pricing power becomes \emph{surplus-claim}
unless it is actively constrained---by anti-rent rules, by exit, or by public
accountability. Each link is individually familiar from the study of rent; the claim
is only that they compose, and that low polycentricity supplies the first link.

A fuller treatment must defend each link empirically and confront the counter-cases---
a centralised public authority that distributes surplus through taxation; a
polycentric system that concentrates it through local oligarchy or validator
cartels---and that empirical defence is the principal work we leave to a successor
paper (\S\ref{sec:limits}).

The relation is directional and probabilistic, not deterministic. High P does not
\emph{guarantee} distributed S---a polycentric system can still tolerate rent---but
low P reliably \emph{erodes} it. This asymmetry is exactly why the framework treats P
and S as independent \emph{measurements} while positing a tendency between them: one
must measure S directly, because P does not determine it; but a low P score is a
leading indicator that S is under structural pressure. The design corollary is that
distributing power is necessary but not sufficient for distributing surplus---the
anti-rent rules must be written explicitly, not assumed to follow.

% ─────────────────────────────────────────────────────────────────────────────
\section{Applying the Framework: Architectural Ideal Types}
\label{sec:application}

We apply the framework to a set of \emph{architectural ideal types}---abstract design
configurations, named by neutral architectural descriptors rather than ideological
labels, and distinct from both deployed systems and the design proposal that follows.
This separation is deliberate: scoring \emph{real} systems is the empirical programme
we defer, and comparing a preferred design against caricatures would stack the deck.
Table~\ref{tab:types} reports illustrative placements.

\begin{table}[!ht]
\centering\small
\renewcommand{\arraystretch}{1.2}
\begin{tabularx}{\textwidth}{@{}Xccc l@{}}
\toprule
\textbf{Architectural ideal type} & \textbf{MCB} & \textbf{P} & \textbf{S} & \textbf{$\Iax$} \\
\midrule
Identity-linked programmable CBDC (state-administered) & 2.0 & 1.0 & 2.0 & low \\
Closed custodial platform \emph{(worked, \S\ref{sec:scoring})} & 1.5 & 0.5 & 0.5 & very low \\
Market-sovereign, no-demurrage design & 1.0 & 3.0 & 1.0 & mid \\
Market-forward convertible design, anti-rent funding & 3.0 & 4.0 & 4.5 & mid \\
Commons-forward demurrage design & 4.5 & 3.0 & 4.5 & mid--high \\
\bottomrule
\end{tabularx}
\caption{Illustrative MCB/P/S placements for five architectural ideal types (half-point
scale), with $\Iax$ as a separate, coarse qualitative band. \emph{Only the closed
custodial platform is worked in full} (Table~\ref{tab:worked}); the remaining
placements are \emph{schematic}---they indicate relative position in the framework and
would require published indicator tables before use as measurements. The lowest
$(P,S)$ placements---programmable CBDC and closed custodial platform---illustrate the
$P \rightsquigarrow S$ tendency: concentrated decision rights co-occurring with
concentrated surplus-claim, because centralised issuance and custody place both the
decision and the residual claim in the same hands.}
\label{tab:types}
\end{table}

Two readings stand out, both about the \emph{shape} the framework reveals rather than
any specific number. First, the configurations that concentrate decision rights (low
P) are the same ones that concentrate surplus-claim (low S)---the $P \rightsquigarrow
S$ tendency made visible. Second, no configuration dominates on all axes: a
market-forward design may score well on P and S while extending a narrower moral
circle than a commons-forward one. The framework surfaces these trade-offs rather
than collapsing them into a ranking.

\subsection{A Hypothesis: the Federated ``Monetary Ecology''}

The configurations that score broadly share a structural form: they are \emph{plural
and layered} rather than monolithic. We therefore advance a \emph{hypothesis}, stated
as such to avoid the circularity of grading a design against axes chosen to favour it:

\begin{quote}
\emph{Hypothesis.} A recursive, federated architecture---minimal federation
invariants at the base (exit rights, proofs of reserves, common formats); a neutral
settlement and reserve layer; a privacy-preserving community circulation layer under
rotating multi-party custody; and a bottom-up mutual-credit layer---tuned per
community through published parameters, \emph{would, if implemented with these
constraints, score highly} on MCB/P/S/$\Iax$.
\end{quote}

The hypothesis is plausible by construction (local, contestable, recallable decision
rights give high P; explicit anti-rent rules support high S; federated custody and
public proofs give high $\Iax$; communities may extend standing per their charters,
admitting broad MCB), but plausibility by construction is not empirical demonstration:
whether such an architecture achieves these scores in deployment, and how it scores
against real systems graded by the same indicators, this paper does not settle. The
architecture's full design and an implementation path are developed in the originating
essay; here it is the framework's leading \emph{implication}, not its centre.

The structural feature that makes such plurality corrigible rather than merely diverse
is \emph{exit}: because participants and communities can leave a configuration they
reject, taking value and state with them, no layer can hold them against the threat of
departure---Hirschman's~\cite{hirschman1970} exit, doing what voice alone cannot.

% ─────────────────────────────────────────────────────────────────────────────
\section{Scope and Limitations}
\label{sec:limits}

\begin{enumerate}
\item \emph{Empirical scoring deferred.} Scoring deployed systems---fiat, mobile
  money, Bitcoin, custodial and stablecoin platforms, CBDC pilots, mutual-credit
  networks---against published indicator evidence is the framework's strongest
  validation and a programme in its own right; it is future work.
\item \emph{$P \rightsquigarrow S$ awaits empirical defence.} \S\ref{sec:causality}
  gives the causal chain; what remains is to defend each link empirically and confront
  the counter-cases---centralised public authorities that distribute surplus,
  polycentric systems that concentrate it through local oligarchy or cartels. This is
  the principal theoretical work the framework still owes.
\item \emph{MCB is the normative axis.} MCB is, by nature, the most normative and
  least mechanical dimension; the framework measures the standing a design encodes,
  not which standings are right.
\item \emph{$\Iax$'s hard cases are open.} The boundary of ``people'' versus
  ``power'' (small co-operatives, validators, municipal treasuries, large merchants),
  and proving reserves and flows without exposing individual transaction graphs, are
  where the diagnostic could become technically distinctive; this paper poses them
  rather than resolving them.
\item \emph{Comparative, not predictive.} The framework renders structural
  commitments legible; it does not forecast adoption. Its conceptual spine is drawn
  from the literatures on money as a social relation~\cite{ingham2004}, polycentric
  governance~\cite{ostrom1990}, collective action and rent~\cite{olson1965},
  exit~\cite{hirschman1970}, and platform concentration~\cite{srnicek2017}; a fuller
  engagement with monetary economics specifically---endogenous money, seigniorage,
  CBDC governance, custodial and payment-rail risk, financial inclusion and
  surveillance---is sized for the empirical successor paper.
\end{enumerate}

% ─────────────────────────────────────────────────────────────────────────────
\section{A Note on Broader Significance: Substrate Design}
\label{sec:broader}

The framework was developed for monetary architecture, but its structure is not
specific to money, and we record the connection briefly rather than develop it here.
Companion work argues that coordinator-based architectures fail
structurally~\cite{nicholson2026}, and that coordinator-free, locally-resolved design
is the correct response wherever knowledge is heterogeneous and locally
held~\cite{nicholson2026hlks}. The monetary case maps onto that argument:
Polycentricity is the monetary form of local, contestable resolution; the $P
\rightsquigarrow S$ tendency is the monetary form of the risk that a coordination
chokepoint extracts rent; and $\Iax$ is the monetary form of the legibility-and-exit
requirements by which such a system keeps capture visible and bounded. That the same
prescription---local resolution, plural and contestable power, exit, transparency of
mechanism---arises independently in monetary design and in distributed computing is,
we argue elsewhere~\cite{nicholson2026hlks}, a convergence rather than an analogy.
That broader claim is the companion papers' to make; here it is only a note on why the
monetary framework may be an instance of something more general.

% ─────────────────────────────────────────────────────────────────────────────
\section{Conclusion}

Money is never neutral: a monetary architecture answers, in its rules, who counts,
who decides, and who benefits, and conventional metrics leave those answers implicit.
The MCB/P/S framework makes them explicit and comparable; the $\Iax$ diagnostic
ensures the answers are not claimed through opaque or capturable means; and the $P
\rightsquigarrow S$ tendency names the structural pressure by which concentrated power
becomes concentrated claim. We have demonstrated the method on a worked example and
exercised it across architectural ideal types, deferring the empirical scoring of
deployed systems---and the fuller defence of $P \rightsquigarrow S$---to future work.
The framework does not by itself prescribe a single architecture. It makes the difference
between architectures measurable, which is the precondition for choosing among them on
grounds more honest than technical performance alone.

% ─────────────────────────────────────────────────────────────────────────────
\bibliography{references}

\end{document}
